\documentclass[12pt,a4paper]{article}
\usepackage[utf8]{inputenc}
\usepackage[T2A]{fontenc}
\usepackage[russian]{babel}
\usepackage{amsmath,amssymb}
\usepackage{booktabs,longtable,array}
\usepackage{geometry}
\usepackage{microtype}
\usepackage{xurl}
\geometry{margin=2.3cm}

\title{Аудит покрытия наблюдаемого мира\\ текущей моделью S2T}
\author{}
\date{6 августа 2026 г.}

\begin{document}
\maketitle

\section{Проверяемый вопрос}

Аудит проверяет не красоту первичного образа, а минимальную эмпирическую
полноту: какие надёжно установленные свойства нашего мира текущая версия S2T
не выводит вообще, выводит только условно или воспроизводит с недостаточной
точностью. Входные величины, postdiction и формулы с незамкнутым операторным
механизмом не засчитываются как независимые предсказания.

\section{Численные провалы и условные совпадения}

Наиболее жёсткий независимый тест даёт электрослабая/QCD-ветвь. При сохранении
электромагнитного train-якоря двухпетлевой split-running даёт
\begin{equation}
 M_W=82.0226\,\mathrm{GeV},\quad
 M_Z=92.0616\,\mathrm{GeV},\quad
 \sin^2\theta_W=0.206203,\quad
 \alpha_s(M_Z)=0.080177.
\end{equation}
Соответствующие относительные невязки равны \(+2.06\%\), \(+0.96\%\),
\(-7.62\%\) и \(-32.05\%\). Минимальная threshold-ветвь также закрыта
отрицательно.

Физический CKM red-team даёт ещё более сильный провал:
\begin{equation}
 J_{S2T}=5.1041\cdot10^{-2},
 \qquad
 J_{CKM}=3.12\cdot10^{-5},
\end{equation}
а предсказанный \(\sin\theta_{13}=0.5694\) вместо контрольного \(0.003732\).
Следовательно, наличие CP-инварианта ещё не означает воспроизведения
наблюдаемой иерархической матрицы смешивания.

Напротив, сверхточные строки требуют понижения эпистемического статуса:
\begin{itemize}
 \item \(S_{vac}\) почти совпадает с \(\alpha^{-1}\), но \(\alpha\) является
 train-якорем, а exact-\(\pi^{-4}\) determinant route закрыт отрицательно;
 \item формула \(m_\tau/m_\mu\) точна до \(3.2\cdot10^{-7}\), но её seed и
 проекционная нормировка не выведены;
 \item нейтринное отношение отличается примерно на \(0.90\%\), но число
 \(23\) не получено representation-consistent способом, абсолютная шкала и
 PMNS отсутствуют;
 \item масса Хиггса численно близка, но наследует условную абсолютную шкалу и
 не сопровождается выводом наблюдаемых couplings и ширины.
\end{itemize}

\section{Поправка на коллективный атлас}

Коллективный \(\pi\)-atlas уже даёт ретроспективное численное покрытие
одиннадцати строк:
\begin{equation}
 \alpha_s,\ \sin^2\theta_W,\ 
 \frac{m_b}{m_p},\frac{m_s}{m_p},\frac{m_d}{m_e},\frac{m_u}{m_e},\
 \frac{m_\tau}{m_\mu},\ V_{cb},\
 \Omega_\Lambda,\Omega_{\mathrm{dm}},\Omega_b.
\end{equation}
Для общей базы получено
\begin{equation}
 E(\pi)=0.01827893,\qquad
 x_{\mathrm{best}}=3.16422,\qquad
 E(x_{\mathrm{best}})=0.01587965.
\end{equation}
Поэтому эти наблюдаемые нельзя называть отсутствующими. Их точный статус:
ретроспективно воспроизведены, но не выведены из выбранного оператора,
parent action и правила адресации observable к формуле.

\section{Минимальный реестр невыведенного}

С учётом атласа текущая версия всё ещё не выводит как единое следствие
первичного образа:
\begin{enumerate}
 \item общий mass/Yukawa-оператор, абсолютные шкалы и running полного
 фермионного спектра;
 \item четыре совместных параметра CKM и полную PMNS-матрицу, несмотря на
 отдельную формулу для \(V_{cb}\);
 \item одновременно \(G_F,M_W,M_Z,\sin^2\theta_W\) и \(\alpha_s(M_Z)\) из
 gauge dynamics;
 \item конфайнмент, адронный спектр, ширины распадов и scattering amplitudes;
 \item единый low-energy action с полным наблюдаемым набором полей,
 аномалий, BRST/BV-комплексом и квантовой мерой;
 \item динамическое пространство-время, нормировку Ньютона и гравитационные
 уравнения;
 \item космологическую динамику расширения, недостающей массы и ускорения,
 несмотря на атласные доли, а также барионную асимметрию.
\end{enumerate}

Отсутствие этих выводов не доказывает ложность первичной гипотезы. Оно
доказывает более узкое утверждение: имеющаяся редукция первичного образа ещё
не идентифицирует наш мир среди эмпирически различимых продолжений.

\section{Вердикт}

На текущем уровне имеются математическое ядро, сильные условные компрессии и
ценные no-go. Однако число закрытых независимых многосекторных физических
предсказаний остаётся равным нулю. Следующий содержательный шаг --- не поиск
ещё одной константы, а вывод общего действия и readout-map, которые без новых
коэффициентов одновременно порождают хотя бы электрослабый/QCD scorecard и
полные матрицы масс и смешивания.

Машинный отчёт сохранён в
\path{s2t_observed_world_coverage_audit.py} и
\path{s2t_observed_world_coverage_results.json}.

\end{document}